\section{Cote 136: an integral associative comparison in every bounded degree}
\label{sec:136associativecomparison}

\paragraph{Question and distinction of structures.}
This concerns the admissible generalized de Rham theories of
cote 136, physical pp.~117 and 119, footers 116 and 118,
\texttt{ARCHIVE\_136\_scan117.tex} and
\texttt{ARCHIVE\_136\_scan119.tex}, and the comparison question
on physical p.~121, footer 120, \texttt{ARCHIVE\_136\_scan121.tex}.
We give the natural derived-section comparison of
AQ-136-COMPARISON an explicit \emph{associative} multiplication.
The intermediate algebras are not asserted to be commutative
or to have compatible divided powers. This is sufficient to
settle the proposed ordinary Massey-product distinction; it
does not settle the archive's strict PD common-source question.

\paragraph{ARCH-136-ASSOCIATIVE-COMPARISON.}
For every simplicial set $X$ there is a nonnegative,
weight-graded associative differential $S=\Gamma_{\mathbb Z}(T)$-
algebra $\mathsf S(X)$, independent of an admissible theory.
For every $C\in\mathcal A_{\mathbb Z}$ there is a natural
zigzag of strict associative differential $S$-algebra maps
\[
 C^{**}(X)\longrightarrow\mathsf K_C(X)
                           \longleftarrow\mathsf S(X).   \tag{136.AC1}
\]
For each $d\geq0$ there is $W_C(d)$ such that, simultaneously
for every $X$ and every $w\geq W_C(d)$, both arrows induce
isomorphisms on $H^j$ for $0\leq j\leq d$ in weight $w$.
Their underlying derived-section maps are the actual natural
comparisons in AQ-136-COMPARISON. No freeness or flatness
of the modules of forms is assumed.

\begin{proof}
\emph{1. A projective resolution with a diagonal.}
Let $U:s\operatorname{Set}\to\operatorname{Set}^{\mathbb N}$
forget all simplicial operators. Its left adjoint is
\[
             F((V_n)_n)=\coprod_{n\geq0}V_n\times\Delta[n].
\]
Write $Q=FU$, with counit $\epsilon$ and comultiplication
$\delta=F\eta U$, where $\eta$ is the unit of the adjunction.
The augmented simplicial object in simplicial sets
\[
 \cdots\ Q^3X\ \substack{\longrightarrow\\[-5pt]
                         \longrightarrow\\[-5pt]
                         \longrightarrow}\ Q^2X
       \ \substack{\longrightarrow\\[-5pt]\longrightarrow}\ QX
           \longrightarrow X
\]
has bar faces $Q^i\epsilon Q^{r-i}$ on $Q^{r+1}X$ and
degeneracies $Q^i\delta Q^{r-i}$. Its simplicial direction
will be called the resolution direction, to distinguish it
from the original simplicial direction of $X$.
Put
\[
 P_r(X)=\mathbb Z[Q^{r+1}X],\qquad
 \partial_r=\sum_{i=0}^r(-1)^i d_i\quad(r>0),\qquad
 \partial_0=0,
\]
using the unnormalized chain complex in the resolution
direction. This is a projective resolution of
$\mathbb Z[X]$ in $s\operatorname{Ab}$.
Indeed, $Q^{r+1}X$ is a coproduct of representable simplicial
sets $\Delta[n]$, and
\[
 \operatorname{Hom}_{s\operatorname{Ab}}
                   (\mathbb Z[\Delta[n]],M)=M_n
\]
is exact. Thus each $P_r(X)$ is projective. After applying
$U$, insertion of the unit $\eta$ gives an extra degeneracy
of the augmentation. Its identities are
$d_0s_{-1}=1$ and $d_i s_{-1}=s_{-1}d_{i-1}$ for $i>0$.
Linearization supplies a contraction in each original
simplicial degree. The augmented complex is therefore exact
degreewise, which is exactness in $s\operatorname{Ab}$.

Every bar term is a free module on a simplicial set, with
its basis-element diagonal. The Alexander--Whitney formula
in the resolution direction defines
\[
 \Delta:P_n(X)\longrightarrow
               \bigoplus_{r+s=n}P_r(X)\otimes P_s(X),
 \qquad
 \Delta[z]=\sum_{r=0}^n
           [z|_{[0,r]}]\otimes[z|_{[r,n]}].             \tag{136.AC2}
\]
Here the two restrictions mean the front and back bar
faces, and tensor products are pointwise in the original
simplicial direction. This is a natural counital,
coassociative chain diagonal. To check coassociativity,
both iterates of (136.AC2) are the sum over two cut
positions $0\leq r\leq s\leq n$ of the same three
successive faces. To check the chain identity, expand
the boundaries of the two successive faces: deleting
an interior vertex agrees with the corresponding cut
of a face of $z$; the two terms deleting the common
cut vertex cancel, with the tensor differential sign.
The endpoints give the counit, which sends each
degree-zero basis element to $1$ and positive resolution
degrees to zero. The augmentation
$P_*(X)\to\mathbb Z[X][0]$ is a coalgebra map: in
degree zero it preserves the basis-element diagonal,
and its other degrees are zero. These are the usual
Alexander--Whitney identities; a primary printed formula
and the identities are Medina-Mardones (2023), \S3.1,
p.~231, equations (3)--(4), DOI 10.4171/EMSS/71.

\emph{2. The convolution algebra and its signs.}
Set
\[
 \mathsf K_C(X)^{w,n}
   =\bigoplus_{r+q=n}
       \operatorname{Hom}_{s\operatorname{Ab}}
                              (P_r(X),C^{w,q}),
 \qquad
 Df=d_hf+(-1)^r d_Cf,                                   \tag{136.AC3}
\]
where $d_hf=f\partial$ when $f$ has resolution degree $r$.
The totalization uses finite diagonals, since $r,q\geq0$.
For $f:P_r\to C^{w,q}$ and $g:P_s\to C^{v,t}$ define
\[
 f\star g=(-1)^{qs}\,
       \mu_C(f\otimes g)\Delta_{r,s}
       :P_{r+s}\longrightarrow C^{w+v,q+t}.             \tag{136.AC4}
\]
Coassociativity of $\Delta$ and associativity of $\mu_C$
give associativity. The signs in the two threefold products
are $qs+(q+t)u$ and $tu+q(s+u)$, which agree modulo two.
The chain-diagonal identity and the Leibniz identity in $C$
give
\[
 D(f\star g)=Df\star g+(-1)^{r+q}f\star Dg.
\]
For explicit signs, the two horizontal terms have exponents
$qs$ and $qs+r$; the two vertical terms have exponents
$r+s+qs$ and $r+s+qs+q$. Expanding the right-hand side
with (136.AC4) gives these same four exponents. The
counit followed by $1\in C^{0,0}$ is the unit.
Constant coefficient cochains define a central map
$S\to\mathsf K_C(X)$: the front or back degree-zero
face in (136.AC2) evaluates a constant coefficient to
that same coefficient, so multiplication on either side
is the ordinary scalar multiplication in $C$.

Let $\underline S$ be the constant simplicial graded
algebra, in form degree zero, and put
$\mathsf S(X)=\mathsf K_{\underline S}(X)$.
The coefficient map gives the right arrow of (136.AC1).
The left arrow is precomposition by
$P_*(X)\to\mathbb Z[X][0]$, using the identity
\[
 C^{w,q}(X)=\operatorname{Hom}_{s\operatorname{Ab}}
                          (\mathbb Z[X],C^{w,q}).
\]
The convolution multiplication on the right of this
identity is exactly pointwise multiplication of sections.
Since the augmentation is a coalgebra map, the left
arrow is multiplicative; it commutes with differential
because the augmentation is a chain map. Both arrows
are therefore strict associative differential $S$-algebra
maps. The entire construction is natural in $X$ and $C$.

\emph{3. The uniform cohomological range.}
The complex (136.AC3) computes
$R\Gamma_X(C^{w,*})$ because $P_*(X)$ is a projective
resolution of $\mathbb Z[X]$. Its maps from ordinary
sections and from $R\Gamma_X(\underline S^w)$ are
precisely those of AQ-136-COMPARISON, realized by this
specific resolution.

For clarity, the two spectral-sequence arguments work
with these explicit terms. In one direction,
projectivity of $P_r$ identifies vertical cohomology
with $\operatorname{Hom}(P_r,H^q(C^{w,*}))$. Uniform
simplexwise vertical admissibility identifies $H^0$
with $\underline S^w$ and kills positive $H^q$ through
the required finite range. In the other direction,
the resolution cohomology of
$\operatorname{Hom}(P_*,C^{w,q})$ is
$R^r\Gamma_X(C^{w,q})$. The complete inherited
AQ-136-COMPARISON proof establishes that a row with
acyclic normalized simplicial complex is
$\Gamma_X$-acyclic for every $X$; the freeness it uses
is that of $N\mathbb Z[X]$, not of the modules of forms.
Thus these groups vanish for $r>0$ once the relevant
row is horizontally acyclic. Take one weight bound
giving vertical and horizontal admissibility through
form degree $d+1$. Both first-quadrant spectral
sequences have finite diagonals and yield the two
asserted cohomology isomorphisms through degree $d$.
All hypotheses used in choosing the bound are uniform
in $X$. This proves the theorem.
\end{proof}

\paragraph{ARCH-136-MASSEY-RIGIDITY.}
For any two $C,D\in\mathcal A_{\mathbb Z}$ there is $N$
such that for every simplicial set $X$, every finite
length $m\geq2$, and all input weights $w_i\geq N$,
the canonical ordinary-cohomology identifications give
equal complete weighted degree-one Massey-product sets
for $C$ and $D$. No connectedness assumption is needed.

\begin{proof}
Take the maximum of $W_C(2)$ and $W_D(2)$ and $1$.
There is a five-term associative algebra zigzag
\[
 C(X)\longrightarrow\mathsf K_C(X)
 \longleftarrow\mathsf S(X)
 \longrightarrow\mathsf K_D(X)\longleftarrow D(X).
                                                               \tag{136.AC5}
\]
In every interval-sum weight at least $N$, each arrow
is an isomorphism on $H^1$ and $H^2$. The complete
triangular gauge-lifting proof of
ARCH-136-MASSEY-TRANSPORT is algebraic: it uses only
associativity, nonnegative cochain degrees,
$H^1$-surjectivity and $H^2$-injectivity in these
weights. It never uses commutativity, divided powers,
or admissibility of the intermediate algebras.
It therefore applies to every arrow of (136.AC5).
The induced identifications are the canonical ones,
by the comparison theorem just proved. Taking the
four equalities proves the assertion, with one bound
independent of $m$ and $X$.
\end{proof}

\paragraph{Remaining gate.}
This closes the ordinary degree-one Massey-product
proposal, including unbounded lengths and primes,
without providing strict arrows in $\mathcal A_{\mathbb Z}$.
The precise distinction is visible in (136.AC2):
the Alexander--Whitney diagonal is generally not
cocommutative, so (136.AC4) need not be commutative.
No divided powers on its convolution algebra have
been constructed. Operations requiring specified
PD relations, and the existence of strict PD common
sources itself, remain outside this proof.
